% Clase 17 --- El solenoide ideal por Ampère (§3.4).
% El punto fuerte no es reobtener B = mu0*n*I, que ya salió de Biot-Savart, sino
% que ahora NO se supuso estar sobre el eje. Por eso la figura lleva DOS curvas:
% la amarilla (un lado adentro, el opuesto afuera) da el valor del campo, y la
% punteada --enteramente adentro, sin corriente encerrada-- prueba que ese valor
% es el mismo a cualquier altura, o sea que el campo es uniforme en la sección.
\begin{center}
\begin{tikzpicture}[scale=1]

  % el solenoide en corte
  \draw[figline, line width=1.1pt] (-4.3,-1.7) -- (4.3,-1.7);
  \draw[figline, line width=1.1pt] (-4.3,1.7) -- (4.3,1.7);
  \foreach \xx in {-4.0,-3.4,-2.8,-2.2,-1.6,-1.0,-0.4,0.2,0.8,1.4,2.0,2.6,3.2,3.8} {
    \node[figblue, scale=0.6] at (\xx,1.7) {$\odot$};
    \node[figblue, scale=0.6] at (\xx,-1.7) {$\otimes$};
  }
  \node[figlblaux, anchor=north east] at (4.3,-1.85) {$n$ espiras por unidad de largo};
  \node[figlblaux, anchor=south east] at (4.3,2.4) {afuera: $\vec B=\vec 0$};

  % campo interior
  \foreach \yy in {-1.2,-0.6,0,0.6,1.2} {
    \draw[figvecres, line width=0.8pt] (-0.7,\yy) -- (0.3,\yy);
  }
  \node[figlbl, figred, anchor=west] at (0.4,1.2) {$\vec B$};

  % ---- curva 1: un lado adentro, el opuesto afuera -> da el valor ----
  \draw[figamber, line width=1.3pt] (-3.6,0.55) rectangle (-1.6,2.75);
  \draw[figvecaux, figamber, line width=1pt] (-2.9,0.55) -- (-2.3,0.55);
  \node[figlbl, figamber, anchor=east, align=right] at (-3.7,0.55)
    {lado interior:\\[-0.25em] $\oint = B\,a$};
  \node[figlbl, figamber, anchor=south, align=center] at (-2.6,2.8)
    {lado exterior: $B=0$};
  \node[figlblaux, figamber, anchor=west] at (-1.5,2.05) {$\vec B\perp d\vec l$};
  \draw[figvecaux] (-2.9,0.2) -- (-3.6,0.2);
  \draw[figvecaux] (-2.3,0.2) -- (-1.6,0.2);
  \node[figlblaux, anchor=north] at (-2.6,0.26) {$a$};

  % ---- curva 2: enteramente adentro -> prueba la uniformidad ----
  \draw[figred, dashed, line width=1.2pt] (1.3,-1.15) rectangle (3.3,1.05);
  \node[figlbl, figred, anchor=west, align=left] at (3.45,-0.05)
    {no encierra corriente:\\[-0.25em] $B$ arriba $=$ $B$ abajo};

  \node[figlbl, anchor=north, align=center] at (0,-2.35)
    {$B\,a = \mu_0 I\,(n\,a) \;\Longrightarrow\; B = \mu_0 n I$,
     \textbf{en cualquier punto interior}};

\end{tikzpicture}\\[0.5em]
{\small\itshape La curva conviene elegirla adaptada a lo que ya se sabe del
campo. Con el rectángulo amarillo ---un lado adentro, paralelo al eje, y el
opuesto afuera--- los dos lados perpendiculares no aportan nada porque ahí
$\vec B\perp d\vec l$, el lado exterior tampoco porque el campo del solenoide
ideal es nulo afuera, y queda sólo $\oint = B\,a$. Del otro lado de la ley, el
rectángulo encierra $n\,a$ espiras llevando $I$ cada una, así que
$B\,a = \mu_0 n a I$ y $B = \mu_0 n I$: el mismo resultado que dio
Biot--Savart, con dos renglones en vez de una integral. Lo importante es lo que
\emph{no} se supuso: en ningún momento hizo falta que el lado interior estuviera
sobre el eje. El rectángulo punteado lo dice de manera aún más directa: al estar
enteramente adentro no encierra corriente, de modo que la circulación es nula y
el campo tiene que valer lo mismo a las dos alturas. Conclusión: dentro del
solenoide ideal $\vec B$ es \textbf{uniforme en toda la sección}, no sólo sobre
el eje ---siempre lejos de las bocas, donde todo esto deja de valer---.}
\end{center}
